\documentclass[10pt]{article}
\usepackage[T1]{fontenc}
\usepackage[utf8]{inputenc}
\usepackage[french]{babel}
\usepackage{amsmath,amssymb,mathrsfs,mathpazo}
\usepackage{geometry}
\geometry{a4paper,margin=2cm}
\usepackage{graphicx,array,tabularx,tabularray,colortbl,multirow,multicol,booktabs,enumitem,pifont,ulem,fancybox,diagbox,cancel}
\usepackage{pstricks,pstricks-add,pst-plot,pst-tree,pst-node,pst-text,pst-func,pst-math,pst-bspline,pst-3dplot,pst-eucl,pst-solides3d}
\let\pscalculate\undefined
\usepackage{pst-bezier}
\usepackage{tikz,pgfplots,tkz-tab}
\usetikzlibrary{arrows,arrows.meta,patterns,calc,positioning,trees,decorations.markings,shapes.geometric,shapes.misc,shapes.arrows,intersections,decorations.pathmorphing,decorations.pathreplacing,backgrounds,shadows,babel}
\pgfplotsset{compat=1.18}
\usepgfplotslibrary{fillbetween}
\usepackage[np]{numprint}
\usepackage{listings,makecell,pas-tableur,twoopt}
\let\barre\undefined
\usepackage[pstricks]{tablvar}
\definecolor{darkgreen}{RGB}{0,100,0}
\providecommand{\red}{\color{red}}
\providecommand{\blue}{\color{blue}}
\providecommand{\green}{\color{green}}
\providecommand{\black}{\color{black}}
\providecommand{\white}{\color{white}}
\providecommand{\gray}{\color{gray}}
\providecommand{\cyan}{\color{cyan}}
\providecommand{\magenta}{\color{magenta}}
\providecommand{\yellow}{\color{yellow}}
\providecommand{\violet}{\color{violet}}
\providecommand{\orange}{\color{orange}}
\providecommand{\euro}{\texteuro}
\providecommand{\EUR}{\texteuro}
\providecommand{\ieme}{\textsuperscript{e}}
\providecommand{\degres}{\ensuremath{^\circ}}
\providecommand{\degree}{\ensuremath{^\circ}}
\providecommand{\vv}[1]{\overrightarrow{#1}}
\providecommand{\Coord}[1]{#1}
\providecommand{\decosix}{}
\providecommand{\starredbullet}{\ensuremath{\bullet}}
\pagestyle{empty}
\setlength{\parindent}{0pt}
\usepackage{adjustbox}\DeclareUnicodeCharacter{2236}{:}\DeclareUnicodeCharacter{2190}{\ensuremath{\leftarrow}}
\providecommand{\R}{\mathbb{R}}
\providecommand{\N}{\mathbb{N}}
\providecommand{\Z}{\mathbb{Z}}
\providecommand{\Q}{\mathbb{Q}}
\providecommand{\D}{\mathbb{D}}
\providecommand{\C}{\mathbb{C}}
\providecommand{\e}{{\mathrm{e}}}
\providecommand{\vect}[1]{\overrightarrow{#1}}
\providecommand{\vectt}[1]{\overrightarrow{\mathrm{#1}}}
\providecommand{\vvt}[1]{\overrightarrow{\mathrm{#1}}}
\providecommand{\np}[2][]{{#2~#1}}
\providecommand{\vv}[1]{{\overrightarrow{#1}}}
\renewcommand{\d}{{\mathrm{d}}}
\providecommand{\degres}{{^\circ}}
\providecommand{\degre}{{^\circ}}
\providecommand{\blue}{\color{blue}}
\providecommand{\red}{\color{red}}
\providecommand{\enskip}{\quad}
\providecommand{\llbracket}{\lbrack\!\lbrack}
\providecommand{\rrbracket}{\rbrack\!\rbrack}
\providecommand{\euro}{{\text{€}}}
\providecommand{\no}{{\text{n°}}}
\providecommand{\interval}[2]{{[#1;#2]}}
\providecommand{\intervFO}[2]{{[#1;#2[}}
\providecommand{\dotfill}{\ldots}

\providecommand{\ds}{\displaystyle}
\providecommand{\textding}[1]{#1}
\providecommand{\Oij}{\(\mathrm{O};\vec{\imath},\vec{\jmath}\)}
\providecommand{\Oijk}{\(\mathrm{O};\vec{\imath},\vec{\jmath},\vec{k}\)}
\providecommand{\Ouv}{\(\mathrm{O};\vec{u},\vec{v}\)}
\providecommand{\intervFF}[2]{\ensuremath{[#1;#2]}}
\renewcommand{\euro}{\ifmmode\text{\texteuro}\else\texteuro\fi}
\DeclareUnicodeCharacter{00B0}{\ifmmode^{\circ}\else\textdegree\fi}
\DeclareUnicodeCharacter{20AC}{\ifmmode\text{\texteuro}\else\texteuro\fi}
% Énoncés : Éducation nationale. Transcription et corrigés : APMEP, Denis Vergès et contributeurs. Découpage et classement : CoopMaths, Sébastien Lozano, Éric Elter et contributeurs. Adaptation de la présentation : Maths Explorer.
% CC BY-SA : https://coopmaths.fr/www/about/
% https://coopmaths.fr/alea/static/e3c/2021/tex/e3c_2021_01_specimen4_1_cor.tex
\begin{document}
\noindent\textit{APMEP / CoopMaths. CC BY-SA. Présentation adaptée par Maths Explorer.}\par\medskip
\bigskip

\textbf{Question 1}

\medskip

On considère la fonction $f$ définie sur $\R$ par : $f(x)= (x+1)\e^{x}$.

La fonction dérivée $f'$ de $f$ est donnée sur $\R$ par:

\begin{center}
{\renewcommand{\arraystretch}{1.5}
\begin{tabularx}{\linewidth}{|X|X|X|X|}
\hline
\textbf{a.}\quad $f'(x)=\e^{x}$ & $\blue\textbf{b.}\quad f'(x)=(x+2)\e^{x}$ 
& \textbf{c.}\quad $f'(x)=-x\e^{x}$   & \textbf{d.}\quad $f'(0)=0$\\
\hline
\end{tabularx}}
\end{center}

\medskip


\begin{tabular}{@{\hspace*{0.05\linewidth}} | p{0.93\linewidth}}
\textbf{Réponse b.}
\end{tabular}

\bigskip

\textbf{Question 2}

\medskip

Pour tous réels $a$ et $b$, le nombre $\dfrac{\e^{a}}{\e^{-b}}$ est égal à:

\begin{center}
{\renewcommand{\arraystretch}{1}
\begin{tabularx}{\linewidth}{|X|X|X|X|}
\hline
\textbf{a.}\quad $\e^{a-b}$ & \textbf{b.}\quad $\e^{\frac{a}{-b}}$ 
& $\blue \textbf{c.}\quad \dfrac{\e^{b}}{\e^{-a}}$  \rule[-10pt]{0pt}{27pt} & \textbf{d.}\quad $\e^{a}-\e^{-b}$\\
\hline
\end{tabularx}}
\end{center}

\medskip


\begin{tabular}{@{\hspace*{0.05\linewidth}} | p{0.93\linewidth}}
$\e^{a} = \dfrac{1}{\e^{-a}}$ et $\dfrac{1}{\e^{-b}}= \e^{b}$ donc $\dfrac{\e^{a}}{\e^{-b}} = \dfrac{\e^{b}}{\e^{-a}}$

\textbf{Réponse c.}
\end{tabular}

\bigskip

\textbf{Question 3}

\medskip

Soit $\left(u_n\right)$  une suite arithmétique telle que $u_3= \dfrac{9}{2}$ et $u_6=3$.

Alors le premier terme $u_0$ et la raison $R$ de la suite sont :

\begin{center}
{\renewcommand{\arraystretch}{1}
\begin{tabularx}{0.8\linewidth}{|X|X|}
\hline
$\blue \textbf{a.}\quad u_0=6 \text{ et } R=-\dfrac{1}{2}$ \rule[-10pt]{0pt}{27pt}& \textbf{b.}\quad $u_0=\dfrac{1}{2}$ et $R=6$ \\
\hline
\textbf{c.}\quad $u_0=6$ et $R=\dfrac{1}{2}$  \rule[-10pt]{0pt}{27pt} & \textbf{d.}\quad $u_0=\dfrac{3}{2}$ et $R=\dfrac{1}{2}$\\
\hline
\end{tabularx}}
\end{center}

\medskip


\begin{tabular}{@{\hspace*{0.05\linewidth}} | p{0.93\linewidth}}
La suite est arithmétique de raison $R$ donc
$u_6=u_5+R=u_4+2R=u_3+3R$.

On déduit que $3 = \frac{9}{2} + 3R$ donc que $3R = 3 - \frac{9}{2} = -\frac{3}{2}$ soit $R=-\frac{1}{2}$.

$u_3=u_0+3R$ donc $u_0=u_3-3R = \frac{9}{2}+ \frac{3}{2}=\frac{12}{2}=6$

\textbf{Réponse a.}
\end{tabular}

\bigskip

\textbf{Question 4}

\medskip

On considère le programme écrit en langage Python ci-dessous.

\begin{center}
\fbox{\tt
\begin{tabular}{l}
s = 0\\
for i in range(51) :\\
\quad s = s + i
\end{tabular}
}
\end{center}

Quelle est la valeur contenue dans la variable \texttt{s} après exécution du programme ?

\begin{center}
{\renewcommand{\arraystretch}{1.5}
\begin{tabularx}{\linewidth}{|X|X|X|X|}
\hline
\textbf{a.}\quad $51$ & \textbf{b.}\quad $\np{1326}$ 
&$\blue \textbf{c.}\quad \np{1275}$   & \textbf{d.}\quad $\np{2500}$\\
\hline
\end{tabularx}}
\end{center}

\medskip


\begin{tabular}{@{\hspace*{0.05\linewidth}} | p{0.93\linewidth}}
 \og \texttt{for i in range(51)} \fg{} signifie que \texttt{i} prend toutes valeurs entières entre 0 et 50.

La valeur contenue dans la variable \texttt{s} après exécution du programme est donc
$0+1+2+\cdots+50$, soit la somme des 50 premiers entiers de 1 à 50.

D'après le cours, $1+2+\cdots+n = \frac{n(n+1)}{2}$ donc
$1+2+\cdots+50 = \frac{50\times 51}{2}=\np{1275}$.

\textbf{Réponse c.}
\end{tabular}

\bigskip

\textbf{Question 5}

\medskip

La valeur exacte de la somme  $S= 1+ \dfrac{1}{2} + \left (\dfrac{1}{2}\right )^2 + \cdots + \left ( \dfrac{1}{2}\right )^{15}$ est:

\begin{center}
{\renewcommand{\arraystretch}{1.7}
\begin{tabularx}{\linewidth}{|X|X|X|X|}
\hline
\textbf{a.}\quad $\np{1,750030518}$ &$\blue \textbf{b.}\quad 2 - \left ( \frac{1}{2}\right )^{15}$ 
& \textbf{c.}\quad $2 - \left ( \frac{1}{2}\right )^{14}$   & \textbf{d.}\quad $\np{1,999969482}$\\
\hline
\end{tabularx}}
\end{center}

\medskip


\begin{tabular}{@{\hspace*{0.05\linewidth}} | p{0.93\linewidth}}
Soit $(u_n)$ la suite géométrique de premier terme $u_0=1$ et de raison $q=\frac{1}{2}$; donc, pour tout $n$, on a $u_n=u_0\times q^{n} = \left (\frac{1}{2}\right )^n$.

Il s'agit donc de calculer la somme des premiers termes d'une suite géométrique.

$S=u_0 + u_1 + \cdots + u_{15}= u_0\times \dfrac{1-q^{16}}{1-q}
= 1\times \dfrac{1-\left (\frac{1}{2}\right ) ^{16}}{1-\frac{1}{2}}
= 2\times \left ( 1- \left (\frac{1}{2}\right )^{16}\right )
=  2- \left (\frac{1}{2}\right )^{15}$


\textbf{Réponse b.}
\end{tabular}

\vspace{0.5cm}

\end{document}
